% !PRINTED 343
\weberpage{343}

We shall therefore consider here only the triple equations of degree $9$,
and first show how the designation and arrangement of the roots
that we became acquainted with in the previous section for the
points of inflection can be derived from the triple property.

Each of the nine roots occurs in four triples, and in all there
therefore exist $4.9:3=12$ distinct triples. Let us take an
arbitrary one of these triples and denote its
roots by
\[
(0,\,0)\ (1,\,2)\ (2,\,1).
\]
Then let $(1,\,1)$ be an arbitrary fourth root. With the first
three this determines, to begin with, three triples, which we denote by
\[
(1,\,1)\ (0,\,0)\ (2,\,2), \quad (1,\,1)\ (1,\,2)\ (1,\,0), \quad (1,\,1)\ (2,\,1)\ (0,\,1)
\]
and there remain two more roots, which form the
fourth triple with $(1,\,1)$, and which we now denote by
\[
(1,\,1)\ (2,\,0)\ (0,\,2)
\]

We now seek the triple determined by $(0,\,0), (2,\,0)$.
In this triple no roots can occur which occur with either of the
two roots $(0,\,0), (2,\,0)$ in one of the triples already determined, % [sic] German source prints feminine "einer" with neuter "Tripel"; expected "einem"
hence not
\[
(1,\,2)\ (2,\,1)\ (1,\,1)\ (2,\,2)\ (0,\,2),
\]
and for the third root of this triple only $(0,\,1)$
or $(1,\,0)$ remains, and the same two roots are also all that remain
for the triple determined by $(0,\,0)\,(0,\,2)$. But since we can
still interchange $(0,\,2)$ with $(2,\,0)$, we do not restrict
generality if the two further triples
\[
(0,\,0)\ (2,\,0)\ (1,\,0), \quad (0,\,0)\ (0,\,2)\ (0,\,1)
\]
are assumed. We now determine the triples that contain $(1,\,0)$,
of which two are already known. The two others
must contain the roots $(2,\,2), (0,\,1), (0,\,2), (2,\,1)$. But since
$(0,\,1)\,(1,\,0)\,(0,\,2)$ and $(0,\,1)\,(1,\,0)\,(2,\,1)$ are not triples [because
$(0,\,1)\,(0,\,2)$ determines the triple $(0,\,0)\,(0,\,2)\,(0,\,1)$ and $(0,\,1)\,(2,\,1)$ the
triple $(1,\,1)\,(2,\,1)\,(0,\,1)$], we have the further
triples
\[
(2,\,2)\ (0,\,1)\ (1,\,0), \quad (1,\,0)\ (0,\,2)\ (2,\,1).
\]
In the same way we obtain
\[
(0,\,1)\ (2,\,0)\ (1,\,2).
\]
